\documentclass[11pt]{article}
\usepackage[margin=0.88in]{geometry}
\usepackage[T1]{fontenc}
\usepackage{lmodern}
\usepackage{microtype}
\usepackage{xcolor}
\usepackage{hyperref}
\usepackage{booktabs}
\usepackage{array}
\usepackage{amsmath,amssymb,amsthm,mathtools}
\usepackage{enumitem}
\usepackage{fancyhdr}
\usepackage{parskip}
\hypersetup{colorlinks=true,linkcolor=blue!55!black,urlcolor=blue!55!black}
\setlist{nosep,leftmargin=*}
\pagestyle{fancy}
\fancyhf{}
\lhead{Causal Fibres Hypothesis Ladder}
\rhead{Hyperseed Formalization 0014}
\cfoot{\thepage}

\newtheorem{definition}{Definition}[section]
\newtheorem{proposition}[definition]{Proposition}
\newtheorem{conjecture}[definition]{Conjecture}
\newtheorem{criterion}[definition]{Kill criterion}
\newtheorem{remark}[definition]{Remark}
\newcommand{\cH}{\mathcal H}
\newcommand{\cD}{\mathcal D}
\newcommand{\cE}{\mathcal E}
\newcommand{\R}{\mathbb R}
\newcommand{\one}{\mathbf 1}
\newcommand{\supp}{\operatorname{supp}}
\newcommand{\KL}{\operatorname{KL}}
\newcommand{\Fone}{\operatorname{F1}}

\title{Causal Fibres as an Evidence-Gated Hypothesis Ladder:\\
A Hyperseed Formalization of Observability, Representation,\\
Plasticity, Settlement, and Credit}
\author{ProtomegaTron / OpenClaw research formalization}
\date{23 July 2026}

\begin{document}
\maketitle

\begin{abstract}
The causal-fibres programme is best understood not as one architectural claim
but as a sequence of defeasible hypotheses.  This note formalizes its H0--H6
ladder as evidence-gated knowledge formation.  A frozen read/write interface
must first expose a teacher correction; rival representation classes must then
compete at matched budget; selective plasticity and its support estimator must
survive held-out tests; only afterward may predictive-coding settlement and a
learned credit highway claim value.  We define loss- and accuracy-based
teacher-gap closure, factor-specific counterfactual closure, representation
objects, support and drift measures, settlement modes, the amortization gap,
and quantitative kill gates.  The Hyperseed interpretation is methodological:
freezing is a conservative falsifier, the representation contest is pluralistic
hypothesis evaluation, and a symbolic cap is registration into an AtomSpace
rather than post hoc labeling.  Proof sketches show why energy descent does not
identify causal credit, why teacher-clamped settlement is not deployment
performance, and why orthogonal fibres are sufficient under common reducing
subspaces but are not a universal model of neural representation.
\end{abstract}

\setcounter{tocdepth}{1}
\tableofcontents
\newpage

\section{Scope, provenance, and epistemic status}

The primary sources are the \texttt{causal-fibres} 0.4.0 release README,
\texttt{TRANSFORMER\_GPU\_EXPERIMENT.md},
\texttt{REPRESENTATION\_CONTEST.md}, and
\texttt{CRITIQUE\_DRIVEN\_REVISION.md} \cite{cfrelease}.  The intended first
GPU programme uses a six-layer synthetic-grammar student, a GPT-2-small
teacher, and controlled counterfactuals for subject number, object number,
tense, negation, and agreement.  This note formalizes that programme; it does
not report that its GPU gates have passed.

The conceptual novelty is the ordering of claims by what their evidence can
actually establish.  Earlier toy results established conditional mechanics,
including exact preservation under supplied local support, a failure in which
energy fell without correct causal credit (R9), a teacher-clamped/free
deployment gap, underidentification of repeated modules, and recovery of exact
blocks when a common block structure truly exists.  They did not establish
that realistic Transformer states possess a small orthogonal causal
decomposition.  We therefore use the following statuses:
\emph{defined}, \emph{conditionally established}, \emph{empirically passed},
\emph{killed}, and \emph{open}.  In this paper H0--H6 are open empirical
hypotheses.

\section{Experimental objects and evidence records}

\begin{definition}[Frozen interface]
Let a base network with fixed parameters $\theta_0$ expose reads
\[
 h(x)=\bigl(h_{\ell_1}(x),\ldots,h_{\ell_s}(x)\bigr)\in\cH
\]
at selected sites and accept a residual write $r(h(x))$ at specified logits or
hidden states.  The pair $I=(\mathrm{read},\mathrm{write})$ is a
\emph{frozen interface} when $\theta_0$ is held constant while the residual
system is trained.  A controlled co-adaptive interface $I_\rho$ may later
permit rank-$\rho$ base updates inside a stated trust region.
\end{definition}

\begin{definition}[Matched budget]
For a representation arm $R$, define its budget vector
\[
 B(R)=(P_{\rm ever},P_{\rm deploy},a,w,Q,t_{\rm train},
       t_{\rm deploy},\operatorname{rank}r),
\]
where the entries are parameters ever trained, deployment parameters, mean
active features, code width, group count, and wall-clock and output-rank
quantities.  Arms are \emph{matched} when designated coordinates of $B$ are
equal within preregistered tolerances; otherwise all coordinates and a
capacity-normalized comparison must be reported.
\end{definition}

\begin{definition}[Evidence record]
An evidence record for hypothesis $H_i$ is
\[
 e_i=(D_{\rm train},D_{\rm val},D_{\rm held},I,R,B,\sigma,M,C),
\]
where $\sigma$ fixes seeds and software/hardware provenance, $M$ is the vector
of observed metrics, and $C$ identifies controls.  An admission predicate
$G_i(e_i)\in\{0,1\}$ states whether the preregistered gate passed.  Advancement
to $H_{i+1}$ requires $G_i(e_i)=1$; a result from a later mechanism cannot
retroactively supply missing evidence for an earlier gate.
\end{definition}

\section{Teacher-gap closure and factor-specific correction}

Let $L_0,L_R,L_T$ be held-out losses for the frozen base, residual-augmented
model, and teacher, respectively, evaluated on the same examples.  Let
$A_0,A_R,A_T$ denote the corresponding accuracies.

\begin{definition}[Teacher-gap closure]
When $L_0>L_T$, loss closure is
\[
 \Gamma_L(R;I)=\frac{L_0-L_R}{L_0-L_T}.
\]
When $A_T>A_0$, accuracy closure is
\[
 \Gamma_A(R;I)=\frac{A_R-A_0}{A_T-A_0}.
\]
A value $0$ means no closure, $1$ matches the teacher gap, values below zero
degrade the base, and values above one surpass the teacher on the stated
metric.  If the denominator is nonpositive or statistically indistinguishable
from zero, the corresponding closure is undefined and cannot pass a gate.
\end{definition}

Ratios can be unstable near a vanishing denominator.  Every report must
therefore include raw losses or accuracies, denominators, uncertainty
intervals, and per-example paired comparisons in addition to $\Gamma$.

\begin{definition}[Counterfactual factor closure]
For causal factor $j$, let $\tau_j x$ edit only factor $j$ by construction,
and let $\delta_j^M(x)=f_M(\tau_jx)-f_M(x)$ be the model's output edit (for
example, a logit vector).  Define the teacher-referenced edit loss
\[
 E_j(M)=\mathbb E_x\!\left[
   d\bigl(\delta_j^M(x),\delta_j^T(x)\bigr)\right].
\]
Provided $E_j(0)>0$, factor-specific closure is
\[
 \Gamma_j=1-\frac{E_j(R)}{E_j(0)}.
\]
Selectivity is recorded by the off-factor spill
\[
 S_{j\rightarrow -j}=
 \mathbb E_x\sum_{k\ne j}
 d_k\bigl(\delta_j^R(x),\delta_j^T(x)\bigr),
\]
or by a preregistered intervention-confusion matrix.  Large total closure with
poor $\Gamma_j$ or high spill is capacity evidence, not factor alignment.
\end{definition}

\begin{criterion}[H0 interface gate]
For an ordinary backpropagation-trained linear, low-rank, or MLP residual on
the exact frozen reads and writes,
\[
 \max\{\Gamma_L,\Gamma_A\}>0.50
\]
on held-out data, with a confidence interval excluding the preregistered
failure region.  If no reasonable BP capacity curve passes, kill or redesign
the interface before interpreting failures of fibres, PC, or credit routing.
\end{criterion}

\section{Five rival representation classes}

Let $h\in\R^d$ be the frozen read and let $r_b$ be the residual write produced
under representation class $b$.

\begin{definition}[Dense residual]
A dense bottleneck uses an encoder $E:\R^d\to\R^k$ and write matrix
$W:\R^k\to\R^o$:
\[
 z=E(h),\qquad r_b=Wz.
\]
It is a raw-capacity control and makes no modularity claim.
\end{definition}

\begin{definition}[Sparse autoencoder]
An SAE has an overcomplete dictionary $D\in\R^{d\times m}$, $m>d$, with
\[
 h\approx Dz,\qquad z\in\R^m,\qquad \|z\|_0\ll m,
\]
and a steering map $r_b=Wz$.  Sparsity may be imposed by top-$k$, thresholding,
or a stated penalty.  Reconstruction alone does not establish causal meaning.
\end{definition}

\begin{definition}[Grouped sparse dictionary]
Let $\{S_q\}_{q=1}^Q$ be subsets of $\{1,\ldots,m\}$, not necessarily
disjoint.  The code is group sparse under, for example,
\[
 \Omega(z)=\sum_{q=1}^Q\lambda_q\|z_{S_q}\|_2,
\qquad h\approx Dz.
\]
Overlaps permit a dictionary atom to participate in more than one candidate
module and hence model superposition without forcing orthogonality.
\end{definition}

\begin{definition}[Orthogonal fibres]
An orthogonal fibre model posits
\[
 \cH=U_1\oplus\cdots\oplus U_Q,\qquad
 U_p\perp U_q\ (p\ne q),
\]
with orthogonal projections $P_q$, code $z_q=P_qh$, and
$r_b=\sum_q W_qz_q$.  A causal interpretation additionally requires selective
interventions and identification; a low off-block operator loss alone is only
structural evidence.
\end{definition}

\begin{definition}[Contextual fibre bundle]
Let $c\in\mathcal C$ index context, $z_{\rm canonical}$ be a shared coordinate,
and $T_c$ an invertible or controlled low-rank local transport:
\[
 z_{\rm local}=T_cz_{\rm canonical},\qquad
 r_b=W_cz_{\rm local}.
\]
Transition maps $T_{c'c}=T_{c'}T_c^{-1}$ relate charts.  This class retains a
canonical intervention or symbolic address while allowing context-dependent
rotations.
\end{definition}

\begin{criterion}[H1 representation gate]
At matched budget, a structured arm must either (i) strictly improve a
preregistered primary metric relative to both dense and SAE-plus-steering
controls, with multiplicity-aware uncertainty, or (ii) add complementary value
on intervention selectivity, held-out combinations, or continual drift while
remaining noninferior on task performance.  Otherwise the structured claim is
killed, even if reconstruction is good.
\end{criterion}

\section{The H0--H6 ladder}

\begin{definition}[H0: frozen interface observability]
There exists a reasonable BP residual class on frozen interface $I$ that closes
more than half the teacher gap, including preregistered factor corrections.
H0 asks whether the correction is visible and writable at all; it does not ask
which structured ontology is correct.
\end{definition}

\begin{definition}[H1: representation contest]
Among matched-budget rivals, at least one structured representation---SAE,
grouped sparse dictionary, orthogonal fibres, or contextual bundle---beats or
complements dense and SAE controls on correction plus an explicitly structural
criterion.  H1 may select different winners in different regimes.
\end{definition}

\begin{definition}[H2: selective persistent plasticity]
Let parameters be partitioned into candidate supports
$\Theta=\bigcup_q\Theta_q$.  For task/context $c$, let $s_c\in[0,1]^Q$ be
supplied or learned responsibility and update
\[
 \Delta\theta_q=-\eta\,s_{c,q}\nabla_{\theta_q}L_c.
\]
For true support $S_c^\star$, define off-support update mass
\[
 D_{\rm off}(c)=
 \frac{\sum_{q\notin S_c^\star}\|\Delta\theta_q\|_2}
      {\sum_q\|\Delta\theta_q\|_2+\epsilon},
\]
and behavioral drift on unrelated probes
\[
 K_{\rm unrelated}(c)=
 \mathbb E_{x\sim D_{\neg c}}\KL
 \bigl(p_{\theta^-}(\cdot|x)\,\|\,p_{\theta^+}(\cdot|x)\bigr).
\]
H2 holds only if gated exact-gradient plasticity reduces preregistered drift
relative to ungated exact gradient, replay, and parameter-isolation controls
while matching current-task performance within tolerance.
\end{definition}

\begin{definition}[H3: held-out generalization]
Train a router or support estimator on contexts $\mathcal C_{\rm train}$ and
evaluate on disjoint contexts or compositions $\mathcal C_{\rm held}$.  With
thresholded prediction $\widehat S_c$,
\[
 P=\frac{|\widehat S_c\cap S_c^\star|}{|\widehat S_c|},\quad
 R=\frac{|\widehat S_c\cap S_c^\star|}{|S_c^\star|},\quad
 \Fone_c=\frac{2PR}{P+R}.
\]
H3 requires mean held-out support $\Fone>0.70$ and generalization of
intervention fingerprints
\[
 \phi_{qj}(c)=
 d_j\!\left(f(x;c),f(\operatorname{do}(z_q\leftarrow z_q');c)\right)
\]
to unseen contexts, not merely in-distribution support recovery.
\end{definition}

\begin{criterion}[H3 prerequisite gate]
If held-out support $\Fone\le0.70$, do not advance to PC/ePC or learned
highways.  A supplied oracle mask may still test conditional mechanics but
must be labeled oracle and cannot count as learned generalization.
\end{criterion}

\begin{definition}[H4: PC/ePC settlement value]
Let $g_{\rm ff}$ be a feed-forward predictor, $g_{\rm rec}$ a matched recurrent
control, and $S_K$ a $K$-step PC/ePC settlement operator.  At matched deployment
compute and parameters, H4 holds only if $S_K$ improves a preregistered
quantity---sample efficiency, robustness, retention, representation quality,
or loss---beyond both controls.  Energy reduction by itself is not such a
quantity.
\end{definition}

\begin{definition}[Settlement modes and amortization gap]
For a common test distribution, distinguish
\[
\begin{array}{ll}
y_{\rm pred}=g_{\rm ff}(h),&
y_{\rm free}=S_K(h;\varnothing),\\
y_{\rm con}=S_K(h;\xi_{\rm test}),&
y_{\rm clamp}=S_K(h;y_T),
\end{array}
\]
where $\xi_{\rm test}$ is admissible retrieval, memory, recurrence, or symbolic
constraint information available at deployment, while $y_T$ is an oracle
teacher signal.  For a loss $\mathcal L$, define
\[
 \Delta_{\rm amort}
 =L_{\rm free}-L_{\rm clamp},\qquad
 L_m=\mathbb E\,\mathcal L(y_m,y^\star).
\]
One may also report the normalized gap
$\widetilde\Delta_{\rm amort}=
(L_{\rm free}-L_{\rm clamp})/(L_{\rm pred}-L_{\rm clamp})$ when its denominator
is positive.
\end{definition}

\begin{criterion}[H4 deployment kill metric]
Predictor-only, free, constrained, and teacher-clamped results must be reported
separately.  A large positive $\Delta_{\rm amort}$ with no matched-compute
free-settlement improvement kills a deployment claim for PC, even if clamped
settlement is excellent.  The clamped result remains an oracle diagnostic.
\end{criterion}

\begin{definition}[H5: constrained settlement exploitation]
H5 holds when admissible $\xi_{\rm test}$ contains genuinely new test-time
information and constrained settlement exploits it better than controls given
the same information and compute:
\[
 L_{\rm con}<\min(L_{\rm free},L_{\rm ff+\xi},L_{\rm rec+\xi})-\delta.
\]
The information source, access policy, latency, and control exposure must be
recorded.  Repackaging training-time teacher labels as a ``constraint'' does
not satisfy H5.
\end{definition}

\begin{definition}[H6: learned causal highway]
Let $g_{\rm exact,c}$ be exact autograd credit after responsibility gating and
$\widehat g_{\psi,c}$ a learned credit highway.  Define directional error
\[
 e_{\rm dir}=1-
 \frac{\langle\widehat g_{\psi,c},g_{\rm exact,c}\rangle}
 {\|\widehat g_{\psi,c}\|\,\|g_{\rm exact,c}\|+\epsilon},
\]
local communication cost $C_{\rm comm}$, update cost $C_{\rm upd}$, and
held-out post-update loss $L_{\rm post}^{\rm held}$.  H6 holds only if the
highway improves at least one preregistered target---locality, novel-context
transfer, or total cost---while remaining noninferior in task and drift metrics
to exact gated credit.  Beating dense feedback alone is insufficient.
\end{definition}

\section{Proof sketches and limits}

\begin{proposition}[Energy descent does not identify causal credit]
A state or parameter update can decrease a scalar energy and improve the
current task while assigning credit to the wrong causal support.
\end{proposition}
\begin{proof}[Proof sketch]
Take two coordinates $(u,v)$ and task output $y=u+v$ with energy
$E(u,v)=\tfrac12(u+v-1)^2$.  Suppose the intervention semantics designate
$u$ as the only cause that should persistently change, while $v$ controls an
unrelated capability.  At $(0,0)$, updates $(\eta,0)$ and $(0,\eta)$ for
$0<\eta<2$ yield exactly the same lower energy
$\tfrac12(\eta-1)^2$.  The second update has zero overlap with the true support
and may damage the unrelated capability, yet energy and current-task loss
cannot distinguish it.  More generally, scalar descent constrains the inner
product with a gradient, not causal identity in a nonidentifiable parameter
class.  The R9 failure is therefore structurally possible: energy descent is
evidence for optimization, not for correct causal credit.
\end{proof}

\begin{proposition}[Teacher-clamped settlement is not deployment performance]
Without a deployment-time channel carrying the teacher constraint, clamped
settlement estimates an oracle conditional risk rather than the free
deployment risk.
\end{proposition}
\begin{proof}[Proof sketch]
Let $H$ be the frozen features and $Y$ the desired output.  A free estimator is
measurable with respect to $\sigma(H)$; a clamped estimator is measurable with
respect to $\sigma(H,Y_T)$, a strictly richer information set whenever
$I(Y;Y_T\mid H)>0$.  By the value-of-information inequality, optimal risk under
the richer sigma-algebra cannot exceed optimal risk under the poorer one.
Thus low $L_{\rm clamp}$ may result entirely from oracle information.  Unless
$Y_T$ or an equivalent constraint is available at deployment, it says nothing
about attainable $L_{\rm free}$.  The observable difference
$\Delta_{\rm amort}$ measures the part the deployable inference process failed
to amortize; it is not repaired by renaming the oracle mode.
\end{proof}

\begin{proposition}[When orthogonal fibres suffice]
For a finite-dimensional real inner-product space $\cH$ and an operator family
$\mathcal A=\{A_c\}$, an orthogonal decomposition
$\cH=\bigoplus_qU_q$ is a common block decomposition exactly when its
orthogonal projections satisfy
\[
 P_qA_c=A_cP_q\qquad\text{for all }q,c.
\]
For self-adjoint operators, common invariant subspaces are reducing and this
condition yields orthogonal fibres.
\end{proposition}
\begin{proof}[Proof sketch]
If $P_q$ commutes with every $A_c$, then $A_cU_q\subseteq U_q$, so all
off-block maps $P_pA_cP_q$ vanish for $p\ne q$.  Conversely, if every $A_c$
preserves every $U_q$ in the orthogonal direct sum, its block matrix has no
off-block entries and commutes with each $P_q$.  For self-adjoint $A_c$,
invariance of $U_q$ implies invariance of $U_q^\perp$, giving a reducing
subspace.
\end{proof}

\begin{remark}[Sufficient, not universal]
The proposition is a conditional sufficiency result, not a representation
theorem for arbitrary neural states.  Nonnormal or context-varying operators
may lack nontrivial common reducing subspaces; noncommuting self-adjoint
operators may generate an irreducible algebra; repeated isomorphic blocks can
make the decomposition nonidentifiable; and superposed features may be better
modeled by an overcomplete dictionary.  Context-dependent conjugacies
$A_c=T_cA_\star T_c^{-1}$ may instead support a bundle whose local charts
rotate.  Orthogonal fibres should therefore win a contest, not be assumed.
\end{remark}

\section{Hyperseed interpretation}

\subsection{The ladder as evidence-gated knowledge formation}

Let $K_i$ be the accepted knowledge state after stage $i$.  The update is
\[
 K_{i+1}=
 \begin{cases}
 K_i\cup\{\langle H_i,e_i,\textsf{supported}\rangle\},&G_i(e_i)=1,\\
 K_i\cup\{\langle H_i,e_i,\textsf{killed}\rangle\},&G_i(e_i)=0.
 \end{cases}
\]
This is a Hyperseed-style formation rule: evidence is not decoration attached
to a favored ontology, but a gate controlling which concepts may acquire
stronger status.  The chain is not logically monotone in the sense that H1
implies H0; rather, it is procedurally conservative: evidence generated under
an unpassed prerequisite cannot license the downstream interpretation.
Alternative surviving representations may coexist with scoped truth values.

\subsection{Freezing as conservative falsifier}

Freezing isolates the proposition ``the correction is readable and writable
through $I$'' from the much larger proposition ``the whole substrate can adapt
until the task works.''  If a frozen BP residual fails H0, adding PC or a
highway only introduces more degrees of freedom and weakens diagnosis.  This
implements a Hyperseed conservativity principle: extend the ontology or
dynamics only after the least expressive admissible account fails under a
recorded test.  Freezing is not dogma; after its result is registered, a
rank-limited $I_\rho$ with anchor drift constraints becomes a new, explicitly
different hypothesis.

\subsection{Pluralistic representation evaluation}

The representation contest instantiates pluralistic hypothesis evaluation.
Each class supplies a different ontology of the same activation:
dense latent capacity, sparse superposed atoms, overlapping functional groups,
orthogonal subspaces, or transported contextual coordinates.  Matched budgets
make these alternatives commensurable without pretending they say the same
thing.  Hyperseed truth values should be indexed by dataset, intervention
family, budget, and context; ``orthogonal fibres supported'' is not promoted
to ``all cognition decomposes orthogonally.''

\subsection{Symbolic cap as registration}

A symbolic cap is scientifically meaningful only when the mapping is
registered before the decisive behavioral test and participates causally in
it.  Let $\alpha_q$ be a neural coordinate or group, $a_q$ an AtomSpace atom,
and $\pi$ provenance.  A minimal MeTTa-style registration is conceptually
\[
\begin{split}
 &(\textsf{CausalCoordinate}\ a_q\ q),\\
 &(\textsf{Grounding}\ a_q\ \alpha_q\ \pi),\\
 &(\textsf{SupportStatus}\ a_q\ \textsf{intervention-aligned}),\\
 &(\textsf{Evidence}\ a_q\ e_i).
\end{split}
\]
A rule-derived atom must then induce a selective neural intervention and the
predicted behavioral consequence on hard negatives and held-out compositions.
Naming a feature after inspecting successful outputs is post hoc
interpretation; it is not registration.  AtomSpace/MeTTa is valuable here
because identity, provenance, evidence status, and rule-mediated effects can
be explicit and queryable \cite{metta,hyperseed}.

\section{Decision table and reporting contract}

\begin{center}
\small
\begin{tabular}{@{}p{0.07\linewidth}p{0.28\linewidth}p{0.29\linewidth}p{0.27\linewidth}@{}}
\toprule
Gate & Claim & Required comparison & Kill consequence\\
\midrule
H0 & Frozen interface exposes correction &
BP residual capacity curve; $\Gamma_L$ or $\Gamma_A>0.50$ &
Change reads/writes or test controlled substrate adaptation.\\
H1 & Structured code adds value &
Dense and SAE+steering at matched budget &
Retain simpler winner; no fibre-emergence claim.\\
H2 & Selective plasticity protects unrelated function &
Ungated exact gradient, replay, isolation &
Do not attribute retention to causal routing.\\
H3 & Support and fingerprints generalize &
Disjoint contexts/compositions; $\Fone>0.70$ &
No PC/highway escalation; oracle masks remain diagnostics.\\
H4 & PC/ePC settlement adds value &
Feed-forward and recurrent controls at matched compute &
Energy descent or clamped success alone is insufficient.\\
H5 & Constraints add deployable information &
All controls receive the same $\xi_{\rm test}$ &
Treat benefit as oracle leakage or ordinary control capacity.\\
H6 & Learned highway beats exact gated credit &
Exact autograd on locality, transfer, cost, task, drift &
Use exact gated credit; do not promote the highway.\\
\bottomrule
\end{tabular}
\end{center}

Every report must include dataset splits, counterfactual generator, frozen and
co-adaptive interfaces, all budget coordinates, raw and normalized task
metrics, factor closure and spill, support precision/recall/F1, off-support
update mass, unrelated KL drift, finite curriculum-order distances, settlement
modes, $\Delta_{\rm amort}$, wall clock, seeds, software commit, and uncertainty.
Negative results remain first-class registered evidence.

\section{Conjectures and open questions}

\begin{conjecture}[Contextual rescue]
When a global orthogonal contest fails but context operators are approximately
conjugate to a stable canonical family, a contextual bundle will improve
held-out intervention fingerprint transport relative to both a single
orthogonal basis and an unconstrained dense adapter at matched budget.
\end{conjecture}

\begin{conjecture}[Selective-plasticity value]
Conditional on H1 and accurate support, responsibility-gated exact gradients
will reduce finite curriculum-order sensitivity and unrelated KL drift without
reducing adaptation on the current task.  Failure would show that the selected
representation is not a useful persistent-plasticity unit even if it predicts
teacher residuals.
\end{conjecture}

Open questions include how uncertainty in teacher quality should propagate
through $\Gamma$; whether support F1 is calibrated to the downstream cost of
false positives versus false negatives; when a bundle's transports are
identifiable rather than gauge choices; and what admissible information policy
distinguishes symbolic constraint from benchmark leakage.  These questions
reinforce the central discipline: later machinery must earn its interpretation
through earlier, scoped evidence.

\section{Conclusion}

The causal-fibres programme becomes sharper when treated as a ladder of
evidence admissions rather than a package of mutually supporting intuitions.
H0 can kill an interface before structure is blamed.  H1 can reject
orthogonality without rejecting useful adaptation.  H2 and H3 distinguish
conditional preservation from learned generalization.  H4 and H5 separate
optimization dynamics from deployable information.  H6 makes a learned
highway compete with the conservative exact-credit baseline.  In Hyperseed
terms, the important object is not ``the fibre'' alone but the fibre together
with its scope, rivals, falsifier, evidence provenance, and permitted
inferences.

\begin{thebibliography}{9}
\bibitem{cfrelease}
OpenClaw/OmegaClaw research workspace,
\emph{causal-fibres 0.4.0 release}: README, Critique-Driven Revision,
Representation Contest, and First GPU Programme, July 2026.

\bibitem{hyperseed}
B. Goertzel,
\emph{Hyperseed v2 source bundle and ontology writings},
preserved in the OpenClaw research library, accessed July 2026.

\bibitem{metta}
SingularityNET,
\emph{MeTTa and Hyperon documentation},
\url{https://metta-lang.dev/}, accessed July 2026.

\bibitem{pc}
R. P. N. Rao and D. H. Ballard,
``Predictive coding in the visual cortex: a functional interpretation of some
extra-classical receptive-field effects,''
\emph{Nature Neuroscience}, 2(1):79--87, 1999.

\bibitem{sae}
T. Bricken et al.,
``Towards Monosemanticity: Decomposing Language Models With Dictionary
Learning,'' Transformer Circuits Thread, 2023.
\end{thebibliography}

\end{document}
